\section{历史财务：增长与资本占用同时观察}

德赛西威过去三年的收入和归母利润持续增长，2025年经营现金流改善更快。与普通消费电子不同，汽车电子项目具有开发期、量产爬坡、客户验收和回款周期，资产负债表与现金流因此是检验收入质量的必需维度。2025年审计年报是本章的事实基础；2026Q1仅用于趋势观察，未被年化。

\begin{exhibitbox}[财务交付与营运资本]
\small
\begin{tabularx}{\textwidth}{L{2cm}Y Y Y Y Y}
\toprule
期间 & 收入（CNYbn） & 归母（CNYbn） & 毛利率 & 经营现金流（CNYbn） & 应收/存货（CNYbn） \\
\midrule
2023A & 21.908 & 1.547 & 19.97\% & 1.141 & 7.168 / 3.260 \\
2024A & 27.618 & 2.005 & 19.88\% & 1.494 & 9.604 / 3.697 \\
2025A & 32.557 & 2.454 & 19.07\% & 2.884 & 9.779 / 4.789 \\
2026Q1 & 6.495 & 0.461 & 18.60\% & 1.116 & 6.638 / 5.492 \\
\bottomrule
\end{tabularx}
\sourcenote{公司2025年年报和2026Q1报告；2026Q1未经审计且不年化。}
\end{exhibitbox}

2025年收入同比增长17.88\%，归母同比增长22.38\%，经营现金流同比增长93.09\%。这组数据支持公司已从单一产品供应走向更大规模的系统交付。另一方面，应收、库存和合同负债的绝对规模也要求读者持续检验现金转换。2026Q1期末应收较年末下降、合同负债上升，但一个季度不能证明全年经营趋势。

\section{分部结构：座舱是基底，智驾是增长，但不是高毛利保证}

座舱占2025年收入63.23\%，是最大的盈利基础；其18.83\%毛利率高于智能驾驶。智能驾驶占29.79\%，增速32.63\%，但毛利率16.36\%、同比下降3.55个百分点。两项事实共同决定我们的建模方式：可以给予智驾已确认收入的增长权重，却不能给予未披露订单、ASP或软件化溢价的额外利润信用。

\begin{exhibitbox}[从收入到EPS的合并口径桥]
\small
\begin{tabularx}{\textwidth}{L{2.5cm}Y Y Y}
\toprule
项 & 2025A & 2026E基准 & 建模处理 \\
\midrule
智能座舱收入 & CNY20.585bn & CNY23.20bn & 已确认收入基础，假设温和增长；不拆分客户/车型ASP \\
智能驾驶收入 & CNY9.700bn & CNY12.10bn & 维持高于整体的增长，但不假定毛利率上升或客户份额扩张 \\
网联及其他收入 & CNY2.272bn & CNY2.70bn & 温和增长代理；不按纯软件给估值 \\
合并收入 & CNY32.557bn & CNY38.00bn & 由产品线相加；外部CNY38.40bn仅为交叉检查 \\
归母/EPS & CNY2.454bn / 已实现 & CNY2.80bn / CNY4.69 & 合并情景，不分配给未披露的产品级净利 \\
\bottomrule
\end{tabularx}
\sourcenote{公司年报、公开券商盈利预测、本院保守合并情景。}
\end{exhibitbox}

本院基准收入同比约增16.7\%，归母约增14.1\%，略低于部分公开券商的CNY38.4bn/CNY2.914bn预测。保守差额不是对公司能力的否定，而是对四项尚未披露字段的折现：客户/车型份额、ASP、订单转收入和智驾增量费用。模型也不把2026Q1低增直接年化为全年熊市，因为季度季节性、项目交付与确认节奏尚未闭合。

\section{研发、全球制造与资产负债表的双面性}

2025年研发投入CNY2.637bn，占收入8.10\%。研发网络、功能安全、算法、测试能力和海外制造网络是系统交付的必要投入，也解释了公司与单一零部件供应商的差异。西班牙工厂、墨西哥量产及国内产能布局可以改善交付半径和客户服务，但其设计产能、利用率、项目毛利和单厂现金回报没有披露。

投资含义是：研发和资本支出应被视为维持产品迭代与全球交付的成本条件，而不是自动转化为收入的新增资产。后续应同时跟踪研发费用率、固定资产投入、库存、应收、减值和经营现金流。若产品收入增长不能被现金验证，成长估值应当收缩。

\section{下一季财务检查表}

\begin{exhibitbox}[中报必须回答的经营问题]
\small
\begin{tabularx}{\textwidth}{L{3cm}Y Y}
\toprule
检查项目 & 支持基础情景的表现 & 降级信号 \\
\midrule
收入 & 重新回到中双位数增长并能解释座舱/智驾贡献 & 连续低个位数增长或显著低于全年基准 \\
智能驾驶毛利率 & 不低于2025A趋势，且成本/产品结构有解释 & 再次明显下滑，成本快于收入 \\
现金转换 & CFO与利润匹配，应收/存货没有异常快于销售 & 营运资本占用扩大、减值风险上升 \\
新增平台 & 至少出现客户、车型/SOP、收入三者之一的可审计升级 & 仅有定点、合作或订单年化描述 \\
\bottomrule
\end{tabularx}
\sourcenote{本院监测框架。}
\end{exhibitbox}

这张检查表将报告从静态目标价变为可证伪的研究工具。任何一项坏信号不必自动推翻长期能力，但会降低基准情景发生的概率；多项同时出现时，应从市场支持型观察降为高估值风险或观察名单。
